\textbf{Proof.}
Let \(m=m_n=\lfloor\rho n\rfloor\). The partition consisting of the vertex set of \(K_m\) and the \(n-m\) isolated singletons is a clique partition. The clique-partition converse gives
\[
R_2(G_n)\geq\left(\frac mn\right)^2+(n-m)\left(\frac1n\right)^2.
\tag{1}
\]
Since \(m/n\to\rho\in(0,1)\), the right side of (1) converges to \(\rho^2\) and is therefore \(\Theta(1)\).

For comparison, the graph invariants in the cited published DTE bound are explicit:
\[
\alpha(G_n)=n-m+1,
\qquad
d_{\mathrm{avg}}(G_n)=\frac{m(m-1)}n,
\qquad
\lambda(G_n)=m-1.
\tag{2}
\]
Indeed, an independent set contains at most one vertex from the clique and all isolated vertices; the only edges are the \(\binom m2\) clique edges; and the adjacency spectral radius of the disjoint union is that of \(K_m\). Consequently,
\[
\frac1{\alpha(G_n)}=\Theta(n^{-1}),
\qquad
\frac{\sqrt{d_{\mathrm{avg}}(G_n)}}n=\Theta(n^{-1/2}),
\qquad
\frac{\lambda(G_n)+1}{n}=\frac mn=\Theta(1).
\tag{3}
\]
The cited lower display therefore supplies only order \(n^{-1/2}\) on this family, whereas (1) supplies order one. The cited spectral upper display is also order one, so it matches (1) in order under precisely the fixed-\(\rho\), \(q=2\), DTE hypotheses used in (2)--(3). No matching upper conclusion for arbitrary clique partitions or another model class follows from this calculation. \(\square\)